\topsection[sec:conclusion]{Conclusion and future work}

We built an attestation chain for AI pipelines. It links datasets, base models, external models, fine-tuned
models and deployments through signed manifests that embed the digests of their parents. It enforces who may
sign which stage, keeps trust and revocation in a signed, pinned, append-only log, and gates deployment on a
verification of the whole chain. The system is complete from the command line to the dashboard, tested by 168
automated tests and by two CI workflows, and documented phase by phase as the brief requires.

On 48 attack and control variants, every manipulation made after signing was detected, including the bypasses
aimed at the index and at the log itself. The cost is under one per cent of training time, 8.5\,ms per
verification and nothing per prediction. The answer to the first sub-question of the project is therefore yes:
an enterprise can verify cryptographically that an artifact, and everything it descends from, is exactly what
authorised parties attested.

The same suite answers the second sub-question. Poison that is present before signing verifies and deploys,
with a backdoor success rate of 77.2\,\% in the reference run and of 74 to 95\,\% for ten other seeds. Content
checks catch the trigger they know and miss an adaptive one. A backdoored model from a trusted publisher
verifies. Cryptographic attestation is necessary for an AI supply chain and it is not sufficient: it must be
combined with evidence about content, and the two must be reported separately.

\textbf{Future work.} In order of priority: publish the trust log and the release roots to a shared
transparency log; replace long-lived keys by keyless signing and hardware-backed roots; repeat the evaluation
on team hardware and with larger models; add trigger-agnostic detectors as further content checks and measure
them against adaptive triggers; export the provenance graph as a signed bill of materials. The paper is ready
to be adapted to the page limit of the first targeted venue.
